% Generated draft: requires mathematical and source review.
\documentclass[11pt]{article}
\usepackage[margin=25mm]{geometry}
\usepackage{fontspec}
\setmainfont{Latin Modern Roman}
\usepackage{amsmath,amssymb,mathtools,unicode-math}
\setmathfont{Latin Modern Math}
\usepackage{enumitem,needspace,xurl,hyperref}
\setlength{\parindent}{0pt}
\setlength{\parskip}{5pt}
\setlength{\emergencystretch}{3em}
\newcommand{\N}{\mathbb{N}}
\newcommand{\Z}{\mathbb{Z}}
\newcommand{\Q}{\mathbb{Q}}
\newcommand{\R}{\mathbb{R}}
\newcommand{\C}{\mathbb{C}}
\newcommand{\F}{\mathbb{F}}
\newcommand{\A}{\mathbb{A}}
\newcommand{\PP}{\mathbb{P}}
\newcommand{\E}{\mathbb{E}}
\newcommand{\eps}{\varepsilon}
\newcommand{\dd}{\,\mathrm{d}}
\begin{document}
\section*{UnsolvedMath ID: 30004092}
\textbf{Source problem code:} catalogue entry
\subsection{Definitions and mathematical statement}
Let \(K\) be an algebraically closed field equipped with a complete non-Archimedean absolute value, and let \(X\) be a variety over \(K\). Write \(X^{\mathrm{an}}\) for its non-Archimedean analytification. The record does not specify whether the absolute value must be nontrivial, nor does it impose smoothness, properness, or a characteristic restriction.

To make the cohomological notation explicit, under the usual interpretation using real tropical differential forms on the Berkovich analytification, let \(\mathcal A^{p,q}\) denote the sheaf of forms of bidegree \((p,q)\) in the chosen theory. Its Dolbeault differential is
\[
 d''\colon \mathcal A^{p,q}\longrightarrow\mathcal A^{p,q+1},
 \qquad (d'')^2=0.
\]
For integers \(p,q\geq 0\), this convention gives
\[
 H^{p,q}(X^{\mathrm{an}})
 =
 \frac{\ker\!\left(d''\colon\Gamma(X^{\mathrm{an}},\mathcal A^{p,q})\longrightarrow\Gamma(X^{\mathrm{an}},\mathcal A^{p,q+1})\right)}
 {\operatorname{im}\!\left(d''\colon\Gamma(X^{\mathrm{an}},\mathcal A^{p,q-1})\longrightarrow\Gamma(X^{\mathrm{an}},\mathcal A^{p,q})\right)},
 \qquad \mathcal A^{p,-1}=0.
\]
This is a conditional explanation of the notation: the supplied record does not identify the precise sheaf of forms, coefficient convention, or variant of tropical Dolbeault cohomology. In particular, the separately cited weakly smooth theory must not be silently identified with the original theory.

\textbf{Research question.} For varieties \(X\) over fields \(K\) as above, determine the groups \(H^{p,q}(X^{\mathrm{an}})\) for all integers \(p,q\geq0\), and characterize the additional hypotheses under which they are finite-dimensional and satisfy duality statements. Determine also their relationship with the tropicalizations of \(X\). The source specifies neither a duality pairing and support convention nor the tropicalizations and comparison maps intended, so these parts remain research directions requiring further formulation. No universal finiteness or duality assertion is being proposed here.
\subsection{Short English statement}
For a variety over an algebraically closed complete non-Archimedean field, study the tropical Dolbeault cohomology of its analytification, including finiteness, duality, and comparison with tropicalizations. The precise cohomology theory and comparison statements remain unspecified in the supplied record.
\subsection{Sources}
\begin{itemize}
\item Supplied local UnsolvedMath registry record, permanent ID \(30004092\), source problem code \texttt{catalogue entry}, in the Oberwolfach Reports Open Problems collection. Title: \emph{Tropical Dolbeault Cohomology of Non-Archimedean Varieties}. The record gives proposed year 2019 and names no individual proposer.
\item Supplied report citation: \emph{Non-Archimedean Geometry and Applications} (2020), Oberwolfach Reports, pp.~513--575, DOI \texttt{10.4171/owr/2019/8}. \url{https://doi.org/10.4171/owr/2019/8}. The record attributes its problem material to workshop open-problem discussions documented in this report.
\item Supplied literature reference: \emph{Tropical Hodge theory}. \url{https://arxiv.org/abs/1909.12633}.
\item Supplied literature reference: \emph{Weakly smooth tropical cohomology}. \url{https://arxiv.org/abs/2111.05747}.
\item Supplied literature reference: \emph{Oberwolfach report}. \url{https://doi.org/10.4171/owr/2019/8}.
\item Literature assessment dated exactly \texttt{2026-08-22}: the supplied curation labels the problem open, reports substantial computations and curve results, describes general Hodge-number and duality questions as unsettled, and states that the original theory can have infinite-dimensional cohomology and fail Poincar\'e duality. These are supplied dated curation claims, not independently verified assessments.
\item The registry labels the statement \texttt{corrected\_verified} and states that the official report was checked directly. This verification claim is part of the supplied provenance; no independent source consultation or mathematical verification is asserted in this draft.
\end{itemize}
\end{document}
